%========================================
% MatinBook - Stage 02: Font System Test
% Test: Persian, Latin, Math, Code Fonts
% Purpose: Validate all font families render correctly
%========================================

% Add parent directory to search path
\makeatletter
\def\input@path{{../}}
\makeatother

\documentclass{matinbook}

\title{آزمون سیستم فونت‌ها}
\author{تیم توسعه MatinBook}
\date{تابستان ۱۴۰۵}

\begin{document}

%----------------------------------------
% Title Page
%----------------------------------------
\maketitle

%----------------------------------------
% Chapter 1: Persian Font
%----------------------------------------
\chapter{فونت فارسی}

\section{خوانایی و زیبایی}

این یک پاراگراف طولانی فارسی برای ارزیابی خوانایی فونت \lr{XB Niloofar} است.
فونت فارسی باید به گونه‌ای طراحی شده باشد که خواننده بتواند ساعت‌های طولانی
بدون خستگی چشم مطالعه کند. فاصله بین حروف، اندازه نقاط، و شکل منحنی‌ها
همگی در خوانایی متن تأثیرگذار هستند. ما در \lr{MatinBook} فونت پیش‌فرض را
با مقیاس ۱.۲ تنظیم کرده‌ایم تا هماهنگی بهتری با فونت لاتین داشته باشد.

\section{حروف خاص فارسی}

حروف فارسی شامل ۳۲ حرف اصلی است. حروفی که در عربی وجود ندارند:

\begin{itemize}
    \item \textbf{«گ»} --- مانند: گفتگو، بزرگ، برگ
    \item \textbf{«چ»} --- مانند: چرا، هیچ، کاغذ
    \item \textbf{«پ»} --- مانند: پدر، پاسخ، پرونده
    \item \textbf{«ژ»} --- مانند: ژرف، مژده، پژوهش
\end{itemize}

\section{اعداد فارسی}

اعداد فارسی باید در تمام حالت‌ها به درستی نمایش داده شوند:

\begin{itemize}
    \item اعداد ساده: ۰ ۱ ۲ ۳ ۴ ۵ ۶ ۷ ۸ ۹
    \item اعداد دو رقمی: ۱۰ ۲۵ ۴۸ ۷۳ ۹۹
    \item اعداد چند رقمی: ۱۴۰۴ ۲۰۲۶ ۱۲۳۴۵۶ ۹۸۷۶۵۴۳۲۱۰
    \item اعداد اعشاری: ۳/۱۴ ۲/۷۱۸ ۱/۶۱۸
    \item درصد: ۲۵٪ ۵۰٪ ۷۵٪ ۱۰۰٪
\end{itemize}

\section{نویسه‌های ویژه}

\begin{itemize}
    \item گیومه فارسی: «نقل قول مستقیم»
    \item گیومه تکی: ‹نقل قول درونی›
    \item پرانتز: (متن داخل پرانتز)
    \item کروشه: [متن داخل کروشه]
    \item آکولاد: \{متن داخل آکولاد\}
    \item ممیز فارسی: ۳/۱۴۱۵۹۲
    \item جداکننده هزارگان: ۱٬۰۰۰٬۰۰۰
\end{itemize}

\section{متن شعرگونه}

\begin{flushright}
    \textbf{توانا بود هر که دانا بود}\\
    \textit{ز دانش دل پیر برنا بود}

    \vspace{0.5cm}
    به نام خداوند جان و خرد\\
    کزین برتر اندیشه برنگذرد
\end{flushright}

%----------------------------------------
% Chapter 2: Latin Font
%----------------------------------------
\chapter{فونت لاتین}

\section{\lr{Times New Roman}}

\begin{latin}
    This chapter demonstrates the Latin font \textbf{Times New Roman},
    which is the default Latin font in MatinBook. This serif font is
    widely used in academic and technical publications due to its
    excellent readability and professional appearance.

    \subsection*{Pangrams (All Letters Test)}

    The quick brown fox jumps over the lazy dog.

    Pack my box with five dozen liquor jugs.

    How vexingly quick daft zebras jump!

    \subsection*{Ligatures Test}

    Common ligatures in English text:

    \begin{itemize}
        \item fi --- finance, define, profile
        \item fl --- flower, reflect, inflate
        \item ff --- offer, different, difficult
        \item ffi --- efficient, coefficient
        \item ffl --- shuffle, waffle
    \end{itemize}

    \subsection*{Special Characters}

    \begin{itemize}
        \item Quotes: ``double'' `single'
        \item Dashes: -- (en-dash), --- (em-dash)
        \item Symbols: @ \# \$ \% \^{} \& * \_ + = | \~{} < > /
        \item Brackets: [square] \{curly\} (round)
    \end{itemize}
\end{latin}

\section{متن علمی انگلیسی}

\begin{latin}
    \subsection*{Computer Science Terms}

    The time complexity of the algorithm is \(O(n \log n)\) in the
    average case and \(O(n^2)\) in the worst case. We use a
    \texttt{HashMap} data structure for \(O(1)\) lookup time.

    \subsection*{Mathematics Terms}

    The Gaussian integral is defined as:
    \[\int_{-\infty}^{\infty} e^{-x^2}\,dx = \sqrt{\pi}\]

    The Riemann zeta function for \(s > 1\):
    \[\zeta(s) = \sum_{n=1}^{\infty} \frac{1}{n^s}\]
\end{latin}

%----------------------------------------
% Chapter 3: Math Font
%----------------------------------------
\chapter{فونت ریاضی}

\section{الفبای یونانی}

\begin{center}
\begin{latin}
\begin{tabular}{ccR{3cm}|ccR{3cm}}
\hline
\multicolumn{3}{c|}{\textbf{Uppercase}} & \multicolumn{3}{c}{\textbf{Lowercase}} \\
\hline
\(\Gamma\) & \texttt{\textbackslash Gamma} & Gamma & \(\alpha\) & \texttt{\textbackslash alpha} & alpha \\
\(\Delta\) & \texttt{\textbackslash Delta} & Delta & \(\beta\)  & \texttt{\textbackslash beta}  & beta \\
\(\Theta\) & \texttt{\textbackslash Theta} & Theta & \(\gamma\) & \texttt{\textbackslash gamma} & gamma \\
\(\Lambda\) & \texttt{\textbackslash Lambda} & Lambda & \(\delta\) & \texttt{\textbackslash delta} & delta \\
\(\Pi\)    & \texttt{\textbackslash Pi}    & Pi    & \(\epsilon\) & \texttt{\textbackslash epsilon} & epsilon \\
\(\Sigma\) & \texttt{\textbackslash Sigma} & Sigma & \(\zeta\)  & \texttt{\textbackslash zeta}  & zeta \\
\(\Phi\)   & \texttt{\textbackslash Phi}   & Phi   & \(\theta\) & \texttt{\textbackslash theta} & theta \\
\(\Omega\) & \texttt{\textbackslash Omega} & Omega & \(\lambda\) & \texttt{\textbackslash lambda} & lambda \\
\hline
\end{tabular}
\end{latin}
\end{center}

\section{عملگرها و نمادها}

\begin{itemize}
    \item جمع و ضرب: $\sum_{i=1}^{n} i = \frac{n(n+1)}{2}$,
    $\prod_{k=1}^{n} k = n!$
    \item انتگرال: $\int_{0}^{\infty} e^{-x} \,dx = 1$,
    $\iint_{D} f(x,y) \,dA$
    \item حد: $\lim_{x \to \infty} (1 + \frac{1}{x})^{x} = e$
    \item مشتق: $\frac{d}{dx}(x^{n}) = nx^{n-1}$,
    $\frac{\partial^{2}f}{\partial x^{2}}$
    \item رادیکال: $\sqrt{x^{2}+y^{2}}$,
    $\sqrt[3]{8} = 2$
\end{itemize}

\section{ماتریس و بردار}

\[
A_{3 \times 3} =
\begin{bmatrix}
    a_{11} & a_{12} & a_{13} \\
    a_{21} & a_{22} & a_{23} \\
    a_{31} & a_{32} & a_{33}
\end{bmatrix}, \quad
\vec{v} = \begin{bmatrix} v_{1} \\ v_{2} \\ v_{3} \end{bmatrix}
\]

\[
\det(A) = a_{11}(a_{22}a_{33} - a_{23}a_{32}) - a_{12}(a_{21}a_{33} - a_{23}a_{31}) + a_{13}(a_{21}a_{32} - a_{22}a_{31})
\]

\section{معادلات معروف}

\begin{align}
    e^{i\pi} + 1 &= 0 \label{eq:euler} \\
    \int_{-\infty}^{\infty} e^{-x^{2}} \,dx &= \sqrt{\pi} \label{eq:gaussian} \\
    \sum_{n=1}^{\infty} \frac{1}{n^{2}} &= \frac{\pi^{2}}{6} \label{eq:basel} \\
    \zeta(s) &= \prod_{p \text{ prime}} \frac{1}{1 - p^{-s}} \label{eq:euler-product}
\end{align}

معادله \meqref{eq:euler} زیباترین معادله ریاضی نامیده می‌شود.

%----------------------------------------
% Chapter 4: Code Font
%----------------------------------------
\chapter{فونت کد}

\section{\lr{JetBrains Mono}}

\lr{JetBrains Mono} یک فونت \lr{monospace} مدرن است که مخصوص
برنامه‌نویسی طراحی شده. این فونت دارای لیگچرهای مخصوص کدنویسی است:

\begin{latin}
\begin{minted}{text}
    Comparison:  == != <= >= -> => :: <| |>
    Arrows:      <- -> --> <-- -->>
    Logic:       && || !!
    Math:        ++ -- ** // /= =/=
\end{minted}
\end{latin}

\section{نمونه کد پایتون}

\begin{latin}
\begin{minted}{python}
def greet(name: str, age: int = 0) -> str:
    """Return a greeting message."""
    if age > 0:
        return f"Hello {name}, you are {age} years old!"
    return f"Hello {name}!"

# Test the function
names = ["Ali", "Sara", "Reza"]
for i, name in enumerate(names, start=1):
    print(f"{i}. {greet(name, 20 + i)}")
\end{minted}
\end{latin}
\codecaption{تابع greeting با type hints}
\label{code:greeting}

\section{مقایسه فونت‌های کد}

\begin{latin}
\begin{minted}{text}
    JetBrains Mono:    The quick brown fox  ||  i != 0  ==>  x => y
    Courier New:       The quick brown fox  ||  i != 0  ==>  x => y
\end{minted}
\end{latin}

%----------------------------------------
% Chapter 5: Multilingual Text
%----------------------------------------
\chapter{متن چندزبانه}

\section{ترکیب فارسی و انگلیسی}

این یک پاراگراف ترکیبی است که در آن از واژگان انگلیسی
مانند \lr{function}، \lr{variable}، \lr{array}، \lr{loop}،
\lr{conditional} و \lr{recursion} استفاده شده است.
همچنین ممکن است اعدادی مانند \lr{3.14159} یا \lr{2.71828}
در متن فارسی ظاهر شوند. عبارات ریاضی مانند $f(x) = x^{2}$
نیز در متن فارسی قرار می‌گیرند.

نام دانشمندان خارجی: \lr{Isaac Newton}، \lr{Albert Einstein}،
\lr{Alan Turing}، \lr{Donald Knuth}، \lr{Ada Lovelace}

\section{نقل قول چندزبانه}

\begin{latin}
    \begin{quote}
        ``Any fool can write code that a computer can understand.
        Good programmers write code that humans can understand.''

        \hfill --- \textit{Martin Fowler}
    \end{quote}
\end{latin}

ترجمه: «هر احمقی می‌تواند کدی بنویسد که کامپیوتر بفهمد.
برنامه‌نویسان خوب کدی می‌نویسند که انسان‌ها بفهمند.»

%----------------------------------------
% Summary
%----------------------------------------
\chapter{خلاصه نتایج}

\section{وضعیت فونت‌ها}

\begin{table}[h]
\centering
\caption{وضعیت خانواده فونت‌های MatinBook}
\label{tab:fonts}
\begin{tabular}{|c|c|c|C{5cm}|}
    \hline
    \textbf{نوع} & \textbf{فونت} & \textbf{وضعیت} & \textbf{توضیحات} \\
    \hline
    فارسی & \lr{XB Niloofar} & \textcolor{matingreen}{✅ فعال} & مقیاس ۱.۲، خوانایی بالا \\
    لاتین & \lr{Times New Roman} & \textcolor{matingreen}{✅ فعال} & سریف، استاندارد آکادمیک \\
    ریاضی & \lr{XITS Math} & \textcolor{matingreen}{✅ فعال} & سازگار با \lr{unicode-math} \\
    کد & \lr{JetBrains Mono} & \textcolor{matingreen}{✅ فعال} & لیگچرهای برنامه‌نویسی \\
    \hline
\end{tabular}
\end{table}

\section{نتیجه}

\textbf{تمام فونت‌های \lr{MatinBook} نسخه ۱ با موفقیت آزمایش شدند.}
نمایش فارسی، انگلیسی، ریاضی و کد همگی بدون نقص هستند. 🎉

\end{document}